\section{Discusión}

Los resultados muestran que el álgebra lineal no es solo una forma de escribir
el problema: también explica el comportamiento del sistema. El ejemplo más
claro es el error de orientación. Si la segunda posición se mide con un ángulo
$\alpha+\delta$ en lugar de $\alpha$, cada punto local $\vect{v}$ se desplaza
\begin{equation}
\|R(\alpha+\delta)\vect{v} - R(\alpha)\vect{v}\|
= \|(R(\delta) - I)\vect{v}\|
= 2\left|\sin\tfrac{\delta}{2}\right|\,\|\vect{v}\|,
\end{equation}
donde se usó que $R(\alpha)$ conserva la norma. El error depende de la
distancia del punto al aparato: un grado de error mueve poco un objeto cercano
y mucho uno lejano, como el extremo de una palanca. La curva de la
Figura~\ref{fig:error-alpha} coincide con esta expresión hasta el redondeo
numérico, lo que confirma tanto la teoría como la implementación.

El error de traslación se comporta de otra manera. Un error en $d_x$ o $d_y$
desplaza todos los puntos el mismo vector, sin importar su distancia, porque
la traslación se suma después de la rotación. En consecuencia, en un escáner
de este tipo conviene medir con más cuidado el ángulo que la posición: con los
objetos a unos 80 cm, $1^\circ$ de error equivale a unos 14 mm, mientras que un
error de 14 mm en la cinta métrica es fácil de evitar.

El ruido del sensor tuvo un efecto mucho menor: \cifRuidoMedio{} mm de error
medio frente a \cifErrorCinco{} mm con $5^\circ$ de error en $\alpha$. Esto
indica que, con este sensor, la calidad del mapa fusionado depende más de la
medición de la pose que de la precisión del VL53L0X.

El trabajo tiene limitaciones que conviene precisar. El sensor mide en un solo
plano, por lo que el mapa es bidimensional y un objeto que no corte ese plano
no aparece; la vista 3D de la interfaz es una ayuda visual. El alcance práctico
es de unos 1,2 m y las superficies oscuras o brillantes devuelven menos luz.
El servo SG90 tiene una precisión angular del orden de uno o dos grados, que se
suma al error de orientación ya discutido. Por último, los resultados
cuantitativos provienen de un entorno simulado; el prototipo físico validó el
barrido del servo y su alimentación, y la medición con el sensor real queda
pendiente.
